% =============================================================================
%  cours.tex - Preparation aux entretiens d'algorithmes
%  Cours imprimable couvrant les 21 patterns essentiels + methode + fiche recap.
%
%  Compilation (deux passes pour la table des matieres) :
%     xelatex cours.tex
%     xelatex cours.tex
%  (ou pdflatex cours.tex, deux fois)
% =============================================================================
\documentclass[11pt,a4paper]{article}

\usepackage[utf8]{inputenc}
\usepackage[T1]{fontenc}
\usepackage[french]{babel}
\usepackage[margin=2.2cm]{geometry}
\usepackage[table]{xcolor}
\usepackage{listings}
\usepackage{tcolorbox}
\usepackage{enumitem}
\usepackage{titlesec}
\usepackage{longtable}
\usepackage{array}
\usepackage{fancyhdr}
\usepackage{hyperref}

% --- Couleurs ---------------------------------------------------------------
\definecolor{primary}{HTML}{1F3A5F}
\definecolor{accent}{HTML}{2E86DE}
\definecolor{codebg}{HTML}{F4F7FB}
\definecolor{codekw}{HTML}{1F3A5F}
\definecolor{codestr}{HTML}{27AE60}
\definecolor{codecom}{HTML}{7F8C99}
\definecolor{rulecol}{HTML}{D5DEE8}

% --- Liens ------------------------------------------------------------------
\hypersetup{colorlinks=true, linkcolor=accent, urlcolor=accent, pdftitle={Preparation entretiens - Algorithmes}}

% --- En-tete / pied de page -------------------------------------------------
\pagestyle{fancy}
\fancyhf{}
\fancyhead[L]{\small\color{primary}Preparation entretiens -- Algorithmes}
\fancyhead[R]{\small\color{codecom}\thepage}
\renewcommand{\headrulewidth}{0.4pt}

% --- Titres de section ------------------------------------------------------
\titleformat{\section}{\Large\bfseries\color{primary}}{\thesection.}{0.6em}{}
\titleformat{\subsection}{\large\bfseries\color{accent}}{}{0em}{}

% --- Style du code Python ---------------------------------------------------
\lstdefinestyle{py}{
  language=Python,
  backgroundcolor=\color{codebg},
  basicstyle=\ttfamily\small,
  keywordstyle=\color{codekw}\bfseries,
  stringstyle=\color{codestr},
  commentstyle=\color{codecom}\itshape,
  numbers=none,
  frame=single,
  rulecolor=\color{rulecol},
  breaklines=true,
  showstringspaces=false,
  columns=fullflexible,
  keepspaces=true,
  aboveskip=8pt, belowskip=8pt,
  xleftmargin=6pt, xrightmargin=6pt,
  literate=%
    {é}{{\'e}}1 {è}{{\`e}}1 {ê}{{\^e}}1 {à}{{\`a}}1 {ô}{{\^o}}1
    {î}{{\^i}}1 {ç}{{\c c}}1 {É}{{\'E}}1 {È}{{\`E}}1 {ù}{{\`u}}1,
}
\lstset{style=py}

% --- Boites "Quand l'utiliser" / "Piege" ------------------------------------
\newtcolorbox{infobox}[1]{colback=codebg,colframe=accent,title=#1,
  fonttitle=\bfseries\color{white},boxrule=0.5pt,arc=2pt,left=6pt,right=6pt,top=4pt,bottom=4pt}

\newcommand{\complexite}[2]{%
  \noindent\textbf{Complexite :}\quad Temps \texttt{#1} \quad|\quad Espace \texttt{#2}\par\smallskip}

\newcommand{\leetcode}[1]{\noindent\textbf{Problemes types :} #1\par\medskip}

% =============================================================================
\begin{document}

% --- Page de titre ----------------------------------------------------------
\begin{titlepage}
\centering
\vspace*{3cm}
{\color{primary}\Huge\bfseries Preparation aux Entretiens\par}
\vspace{0.6cm}
{\color{accent}\LARGE Algorithmes \& Structures de donnees\par}
\vspace{1.2cm}
{\large Comment aborder un probleme de code,\\
les 21 patterns essentiels, et les pieges a eviter.\par}
\vfill
{\color{codecom}\large Cours a imprimer -- support d'entrainement\par}
\vspace{2cm}
\end{titlepage}

\tableofcontents
\newpage

% =============================================================================
\section{Methode : comment coder en entretien}
% =============================================================================
Ne jamais foncer sur le clavier. Un entretien evalue autant ton \emph{raisonnement}
que ta solution. Suis systematiquement ces six etapes :

\begin{enumerate}[leftmargin=*,itemsep=4pt]
  \item \textbf{Comprendre} -- reformule l'enonce ; clarifie les contraintes
        (taille, valeurs, doublons, donnees triees ?).
  \item \textbf{Exemples} -- ecris 1 ou 2 exemples a la main, y compris les cas
        limites (tableau vide, un seul element).
  \item \textbf{Force brute} -- enonce la solution naive et sa complexite.
        Tu as toujours une base sur laquelle t'appuyer.
  \item \textbf{Optimiser} -- cherche le pattern : quelle structure de donnees ?
        quel invariant ? un tri prealable aide-t-il ?
  \item \textbf{Coder} -- ecris proprement en parlant a voix haute ; nomme
        clairement les variables.
  \item \textbf{Tester} -- deroule le code sur un exemple, verifie les cas
        limites, corrige.
\end{enumerate}

\begin{infobox}{Regle d'or}
Annonce toujours la complexite (temps \emph{et} espace) de ta solution, et
propose une piste d'amelioration meme si tu n'as pas le temps de l'implementer.
\end{infobox}

\newpage
% =============================================================================
\section{Hash Map / Dictionnaire}
% =============================================================================
Structure cle $\rightarrow$ valeur offrant un acces en \texttt{O(1)} moyen.
C'est l'outil le plus rentable en entretien : il transforme souvent un
\texttt{O(n\textsuperscript{2})} en \texttt{O(n)}.

\textbf{L'idee.} Au lieu de re-parcourir tout le tableau a chaque fois qu'on
cherche quelque chose, on \emph{memorise au fur et a mesure} ce qu'on a deja vu
dans un dictionnaire, puis on le consulte instantanement. On echange de la
memoire (\texttt{O(n)}) contre du temps : c'est le compromis classique
« espace contre vitesse ».

\textbf{Reconnaitre le pattern.} Des que l'enonce parle de \emph{compter}, de
\emph{regrouper}, de \emph{trouver une paire / un complement}, ou que ta premiere
idee est une double boucle qui compare chaque element a tous les autres, pense
hash map : elle supprime presque toujours la boucle interne.

\begin{infobox}{Quand l'utiliser}
Compter des frequences ; retrouver un complement deja vu (Two Sum) ; grouper
des elements par cle ; memoriser des indices deja rencontres.
\end{infobox}

\complexite{O(1) moyen}{O(n)}
\leetcode{Two Sum (LC 1), Group Anagrams (LC 49), Top K Frequent (LC 347).}

\begin{lstlisting}
def two_sum(nums, target):
    seen = {}                      # valeur -> indice
    for i, x in enumerate(nums):
        if target - x in seen:
            return [seen[target - x], i]
        seen[x] = i
    return []
\end{lstlisting}

\textbf{Piege :} en cas de collision de cles, la derniere valeur ecrase la
precedente. Reflechis a ce que la cle doit reellement representer.

% =============================================================================
\section{Hash Set}
% =============================================================================
Ensemble d'elements uniques ; test d'appartenance en \texttt{O(1)} moyen.

\textbf{L'idee.} Un set est un hash map sans valeurs : on ne retient que la
\emph{presence} des elements. Il repond en temps constant a une seule question,
mais la plus utile : « ai-je deja rencontre cet element ? ». On l'alimente au fil
du parcours et on interroge avant d'ajouter.

\textbf{Reconnaitre le pattern.} Mots-cles declencheurs : \emph{doublon},
\emph{unique}, \emph{deja vu}, \emph{visite}. Si tu es tente d'ecrire
\texttt{if x in ma\_liste} a l'interieur d'une boucle (recherche en \texttt{O(n)}
a chaque tour, donc \texttt{O(n\textsuperscript{2})} au total), remplace la liste
par un set et le cout tombe a \texttt{O(n)}.

\begin{infobox}{Quand l'utiliser}
Tester l'appartenance rapidement ; detecter les doublons ; suivre les elements
deja visites (parcours de graphe, cycles).
\end{infobox}

\complexite{O(1) moyen}{O(n)}
\leetcode{Contains Duplicate (LC 217), Longest Consecutive (LC 128), Happy Number (LC 202).}

\begin{lstlisting}
def has_duplicate(nums):
    seen = set()
    for x in nums:
        if x in seen:
            return True
        seen.add(x)
    return False
\end{lstlisting}

\textbf{Piege :} un set ne conserve pas l'ordre d'insertion ; si l'ordre compte,
utilise un dictionnaire (ordonne depuis Python 3.7) ou une liste annexe.

% =============================================================================
\section{Two Pointers (deux pointeurs)}
% =============================================================================
Deux indices qui avancent selon une regle. Deux variantes : \emph{opposes}
(left/right) ou \emph{meme direction} (lent/rapide).

\textbf{L'idee.} Beaucoup de problemes qui semblent demander deux boucles
imbriquees peuvent se resoudre avec deux curseurs qui balaient le tableau en un
seul passage. Dans la variante \emph{opposee}, on part des deux extremites d'un
tableau \emph{trie} et on rapproche les curseurs selon que la somme est trop
grande ou trop petite. Dans la variante \emph{lent/rapide}, un curseur avance
plus vite que l'autre (utile pour les cycles ou le milieu d'une liste).

\textbf{Reconnaitre le pattern.} Tableau ou chaine \emph{trie}, recherche d'une
\emph{paire} qui verifie une condition, \emph{palindrome}, ou l'envie d'ecrire
une double boucle \texttt{O(n\textsuperscript{2})} sur des donnees ordonnees :
les deux pointeurs ramenent souvent le cout a \texttt{O(n)} et \texttt{O(1)}
d'espace.

\begin{infobox}{Quand l'utiliser}
Tableau trie et recherche d'une paire ; palindrome ou inversion ; detection de
cycle et recherche du milieu (lent/rapide).
\end{infobox}

\complexite{O(n)}{O(1)}
\leetcode{Valid Palindrome (LC 125), Two Sum II (LC 167), 3Sum (LC 15).}

\begin{lstlisting}
def two_sum_sorted(nums, target):
    left, right = 0, len(nums) - 1
    while left < right:
        s = nums[left] + nums[right]
        if s == target:
            return [left, right]
        if s < target:
            left += 1
        else:
            right -= 1
    return []
\end{lstlisting}

\textbf{Piege :} la variante opposee suppose un tableau \emph{trie}. Si l'enonce
ne le garantit pas, il faut trier d'abord (\texttt{+ O(n log n)}).

% =============================================================================
\section{Sliding Window (fenetre glissante)}
% =============================================================================
Maintenir une fenetre \texttt{[left, right]} et l'ajuster selon une contrainte,
au lieu de recalculer chaque sous-tableau.

\textbf{L'idee.} Verifier tous les sous-tableaux coute \texttt{O(n\textsuperscript{2})}.
La fenetre glissante evite ce gaspillage : quand on avance le bord droit, on
\emph{ajoute} un element ; quand la contrainte est violee, on avance le bord
gauche pour \emph{retirer} des elements. Chaque element entre et sort de la
fenetre au plus une fois, d'ou un cout total de \texttt{O(n)}. C'est comme une
chenille qui s'etire et se contracte le long du tableau.

\textbf{Reconnaitre le pattern.} L'enonce demande le \emph{plus long} / \emph{plus
court} sous-tableau ou sous-chaine \emph{contigu} respectant une condition (somme,
nombre de caracteres distincts, etc.), ou une fenetre de \emph{taille fixe k}. Le
mot cle est \textbf{contigu} : si les elements doivent etre voisins, pense fenetre
glissante.

\begin{infobox}{Quand l'utiliser}
Sous-tableau ou sous-chaine \emph{contigu} ; longueur min/max sous contrainte ;
fenetre de taille fixe (somme de k elements consecutifs).
\end{infobox}

\complexite{O(n)}{O(1) ou O(k)}
\leetcode{Longest Substring (LC 3), Min Size Subarray (LC 209), Find Anagrams (LC 438).}

\begin{lstlisting}
def longest_unique(s):
    last = {}                      # caractere -> dernier indice vu
    left = best = 0
    for right, ch in enumerate(s):
        if ch in last and last[ch] >= left:
            left = last[ch] + 1
        last[ch] = right
        best = max(best, right - left + 1)
    return best
\end{lstlisting}

\textbf{Piege :} bien decider quand \emph{agrandir} (avancer \texttt{right}) et
quand \emph{retrecir} (avancer \texttt{left}) la fenetre ; un ordre incorrect
casse l'invariant.

% =============================================================================
\section{Prefix Sum (somme prefixee)}
% =============================================================================
On precalcule \texttt{prefix[i] = somme de nums[0..i-1]}. Toute somme
d'intervalle se lit alors en \texttt{O(1)} : \texttt{sum(i, j) = prefix[j+1] - prefix[i]}.

\textbf{L'idee.} Repondre a beaucoup de questions « somme entre i et j » en
re-additionnant a chaque fois coute cher. En calculant \emph{une seule fois} les
sommes cumulees depuis le debut, toute somme partielle devient une simple
soustraction de deux valeurs deja connues. On investit \texttt{O(n)} au depart
pour que chaque requete suivante soit gratuite.

\textbf{Reconnaitre le pattern.} \emph{Nombreuses} requetes de somme (ou de
compte) sur des intervalles, recherche de sous-tableaux dont la somme vaut une
cible, ou notion d'\emph{equilibre} (indice pivot). Tres puissant combine a un
hash map : on stocke les sommes prefixees deja vues pour compter les
sous-tableaux en un seul passage.

\begin{infobox}{Quand l'utiliser}
Nombreuses requetes de somme sur intervalle ; compter les sous-tableaux de somme
donnee (combine avec un hash map) ; indice pivot.
\end{infobox}

\complexite{O(n) precalcul, O(1) requete}{O(n)}
\leetcode{Range Sum Query (LC 303), Subarray Sum = K (LC 560), Pivot Index (LC 724).}

\begin{lstlisting}
def subarray_sum_k(nums, k):
    count = running = 0
    seen = {0: 1}                  # somme prefixee -> nb d'occurrences
    for x in nums:
        running += x
        count += seen.get(running - k, 0)
        seen[running] = seen.get(running, 0) + 1
    return count
\end{lstlisting}

\textbf{Piege :} initialise \texttt{seen = \{0: 1\}} pour compter les
sous-tableaux qui commencent a l'indice 0.

% =============================================================================
\section{Sorting (tri)}
% =============================================================================
Trier revele souvent la structure cachee d'un probleme (voisinage, chevauchement).
Cout de reference : \texttt{O(n log n)}.

\textbf{L'idee.} Une fois les donnees ordonnees, les elements qui « vont
ensemble » deviennent voisins : les doublons se touchent, les intervalles qui se
chevauchent se suivent, une recherche dichotomique devient possible. Trier n'est
pas la reponse en soi, mais l'\emph{etape preparatoire} qui rend le vrai
algorithme simple. En Python, \texttt{sort()} (Timsort) est stable et l'argument
\texttt{key} permet de trier selon n'importe quel critere.

\textbf{Reconnaitre le pattern.} L'ordre des elements n'a pas d'importance dans
l'enonce, mais un traitement devient evident \emph{s'ils sont ordonnes} :
fusionner des intervalles, detecter des doublons, planifier des reunions,
preparer un two-pointers ou une dichotomie. Cout a assumer : \texttt{O(n log n)}.

\begin{infobox}{Quand l'utiliser}
Regrouper ou detecter des doublons ; fusionner des intervalles ; preparer un
two-pointers ; ordonner selon un critere personnalise.
\end{infobox}

\complexite{O(n log n)}{O(n) ou O(1)}
\leetcode{Merge Intervals (LC 56), Meeting Rooms (LC 252), Largest Number (LC 179).}

\begin{lstlisting}
def merge_intervals(intervals):
    intervals.sort(key=lambda p: p[0])
    res = []
    for start, end in intervals:
        if res and start <= res[-1][1]:
            res[-1][1] = max(res[-1][1], end)   # fusionner
        else:
            res.append([start, end])
    return res
\end{lstlisting}

\textbf{Piege :} le tri detruit l'ordre initial. Si tu dois rendre des indices
d'origine, trie des paires \texttt{(valeur, indice)}.

% =============================================================================
\section{Binary Search (recherche dichotomique)}
% =============================================================================
Recherche en \texttt{O(log n)} sur un espace trie ou monotone. Cle : trouver un
\emph{predicat monotone} (vrai puis faux, ou l'inverse).

\textbf{L'idee.} A chaque etape, on regarde l'element du milieu et on elimine
\emph{la moitie} des candidats. Diviser l'espace de recherche par deux a chaque
tour donne un cout logarithmique : chercher dans un million d'elements ne demande
qu'une vingtaine de comparaisons. La version avancee ne cherche pas une valeur
mais la \emph{frontiere} d'une propriete monotone : « quelle est la plus petite
valeur qui fonctionne ? ».

\textbf{Reconnaitre le pattern.} Donnees \emph{triees}, ou reponse qui varie de
facon \emph{monotone} (si \texttt{x} marche, tout \texttt{x} plus grand marche
aussi). Formulation type : « trouver la premiere / derniere position qui\ldots »,
ou « minimiser une capacite / vitesse telle que\ldots ». Des qu'on peut poser une
question a reponse oui/non monotone, la dichotomie s'applique.

\begin{infobox}{Quand l'utiliser}
Tableau trie ; reponse monotone (« la valeur x fonctionne-t-elle ? ») ;
recherche de la premiere ou derniere position valide.
\end{infobox}

\complexite{O(log n)}{O(1)}
\leetcode{Binary Search (LC 704), Search in Rotated Array (LC 33), Koko Bananas (LC 875).}

\begin{lstlisting}
def binary_search(nums, target):
    left, right = 0, len(nums) - 1
    while left <= right:
        mid = left + (right - left) // 2   # evite le debordement
        if nums[mid] == target:
            return mid
        if nums[mid] < target:
            left = mid + 1
        else:
            right = mid - 1
    return -1
\end{lstlisting}

\textbf{Piege :} attention aux bornes (\texttt{<} vs \texttt{<=}) et a la mise a
jour de \texttt{mid $\pm$ 1} pour eviter les boucles infinies.

% =============================================================================
\section{Stack (pile LIFO)}
% =============================================================================
Last In, First Out. En Python : une simple \texttt{list} avec \texttt{append()}
et \texttt{pop()}.

\textbf{L'idee.} La pile retient le dernier element ajoute pour le traiter en
premier -- exactement le comportement dont on a besoin quand des elements doivent
etre « fermes » dans l'ordre inverse de leur ouverture (parentheses, balises),
ou quand on veut suspendre une tache pour y revenir plus tard (DFS iteratif,
annuler/refaire). Les operations \texttt{push} et \texttt{pop} sont en
\texttt{O(1)}.

\textbf{Reconnaitre le pattern.} Correspondances imbriquees (parentheses,
accolades, balises), evaluation d'expressions, parcours en profondeur sans
recursion, ou tout probleme ou « le plus recent d'abord » est la bonne regle.
Signal fort : on doit se souvenir d'elements \emph{en attente} et les traiter dans
l'ordre inverse.

\begin{infobox}{Quand l'utiliser}
Parentheses / expressions equilibrees ; parsing et evaluation ; DFS iteratif ;
fonctionnalites annuler / refaire.
\end{infobox}

\complexite{O(1) par operation}{O(n)}
\leetcode{Valid Parentheses (LC 20), Min Stack (LC 155), Evaluate RPN (LC 150).}

\begin{lstlisting}
def valid_parentheses(s):
    pairs = {")": "(", "]": "[", "}": "{"}
    stack = []
    for c in s:
        if c in pairs:
            if not stack or stack.pop() != pairs[c]:
                return False
        else:
            stack.append(c)
    return not stack
\end{lstlisting}

\textbf{Piege :} pense a verifier que la pile est \emph{vide} a la fin ; une pile
non vide signale des ouvrants non fermes.

% =============================================================================
\section{Monotonic Stack (pile monotone)}
% =============================================================================
Une pile dont les elements restent toujours croissants ou decroissants. Chaque
element est empile et depile au plus une fois : cout total \texttt{O(n)}.

\textbf{L'idee.} Quand on cherche, pour chaque element, « le prochain plus
grand » (ou plus petit), une double boucle donne \texttt{O(n\textsuperscript{2})}.
La pile monotone garde seulement les candidats encore \emph{utiles} : avant
d'empiler un nouvel element, on depile tous ceux qu'il « bat », et c'est
justement lui leur reponse. Comme chaque element entre et sort une seule fois, le
cout total reste \texttt{O(n)} malgre la boucle interne.

\textbf{Reconnaitre le pattern.} Formulations « prochain / precedent element plus
grand / plus petit », spans (combien de jours jusqu'au prochain plus chaud),
histogrammes, eaux de pluie. Astuce : on empile souvent des \emph{indices} plutot
que des valeurs, pour pouvoir calculer des distances.

\begin{infobox}{Quand l'utiliser}
Prochain element plus grand / plus petit ; largest rectangle in histogram ;
trapping rain water ; spans boursiers.
\end{infobox}

\complexite{O(n)}{O(n)}
\leetcode{Daily Temperatures (LC 739), Next Greater Element (LC 496), Largest Rectangle (LC 84).}

\begin{lstlisting}
def daily_temperatures(temps):
    res = [0] * len(temps)
    stack = []                     # indices, temperatures decroissantes
    for i, t in enumerate(temps):
        while stack and temps[stack[-1]] < t:
            j = stack.pop()
            res[j] = i - j
        stack.append(i)
    return res
\end{lstlisting}

\textbf{Piege :} on empile souvent des \emph{indices} plutot que des valeurs,
pour pouvoir calculer des distances.

% =============================================================================
\section{Backtracking (retour sur trace)}
% =============================================================================
Explorer toutes les possibilites via un schema \emph{choisir / recurser /
annuler}. Cout souvent exponentiel.

\textbf{L'idee.} On construit une solution pas a pas. A chaque etape on fait un
\emph{choix}, on explore recursivement ce que ce choix entraine, puis on
\emph{annule} le choix (backtrack) pour essayer le suivant. On explore ainsi un
\emph{arbre de decisions}, et l'annulation permet de reutiliser la meme structure
pour toutes les branches. On peut couper des branches (pruning) des qu'une
solution partielle viole une contrainte, ce qui accelere beaucoup en pratique.

\textbf{Reconnaitre le pattern.} L'enonce demande \emph{toutes} les solutions :
combinaisons, permutations, sous-ensembles, ou un placement respectant des
contraintes (Sudoku, N-Queens). Signal : « generer / enumerer tous les\ldots ».
Attention au cout, naturellement exponentiel (\texttt{O(2\textsuperscript{n})} ou
\texttt{O(n!)}).

\begin{infobox}{Quand l'utiliser}
Toutes les combinaisons, permutations ou sous-ensembles ; problemes a contraintes
(Sudoku, N-Queens) ; generation exhaustive.
\end{infobox}

\complexite{O(2\textsuperscript{n}) ou O(n!)}{O(n) pile de recursion}
\leetcode{Subsets (LC 78), Permutations (LC 46), Combination Sum (LC 39).}

\begin{lstlisting}
def subsets(nums):
    res, path = [], []
    def backtrack(start):
        res.append(path[:])        # copie de l'etat courant
        for i in range(start, len(nums)):
            path.append(nums[i])   # choisir
            backtrack(i + 1)
            path.pop()             # annuler (backtrack)
    backtrack(0)
    return res
\end{lstlisting}

\textbf{Piege :} enregistre une \emph{copie} (\texttt{path[:]}), sinon toutes les
solutions pointeront vers la meme liste mutable.

% =============================================================================
\section{Linked List (liste chainee)}
% =============================================================================
Des noeuds relies par \texttt{.val} et \texttt{.next}. Les entretiens insistent
sur la manipulation propre des pointeurs.

\textbf{L'idee.} Contrairement a un tableau, on ne peut pas sauter directement au
i-eme element : il faut suivre les liens de proche en proche. Tout le savoir-faire
consiste a manipuler quelques pointeurs (\texttt{prev}, \texttt{curr},
\texttt{next}) dans le bon ordre, sans perdre le reste de la liste. Deux
techniques reviennent : le \emph{noeud sentinelle} (dummy) qui simplifie la
gestion de la tete, et les pointeurs \emph{lent/rapide} pour le milieu ou la
detection de cycle.

\textbf{Reconnaitre le pattern.} L'enonce fournit une liste chainee et demande de
l'\emph{inverser}, la \emph{fusionner}, detecter un \emph{cycle}, ou supprimer un
noeud. Le vrai defi n'est pas l'algorithme mais la \emph{rigueur} sur les cas
limites : liste vide, un seul noeud, modification de la tete ou de la queue.

\begin{infobox}{Quand l'utiliser}
Inverser une liste ; detecter un cycle (lent/rapide) ; fusionner deux listes ;
supprimer le n-ieme noeud depuis la fin.
\end{infobox}

\complexite{O(n)}{O(1)}
\leetcode{Reverse List (LC 206), Linked List Cycle (LC 141), Merge Two Lists (LC 21).}

\begin{lstlisting}
def reverse_list(head):
    prev = None
    curr = head
    while curr:
        nxt = curr.next            # sauvegarder le suivant
        curr.next = prev           # inverser le lien
        prev = curr
        curr = nxt
    return prev
\end{lstlisting}

\textbf{Piege :} soigne les cas limites -- liste vide, un seul noeud, mise a jour
de \texttt{head}. Un \emph{noeud sentinelle} (dummy) simplifie souvent le code.

% =============================================================================
\section{Fast \& Slow Pointers (lievre et tortue)}
% =============================================================================
Deux pointeurs avancent a des vitesses differentes : \texttt{slow} d'un pas,
\texttt{fast} de deux. Le rapide finit toujours par rattraper le lent \emph{s'il
existe un cycle}.

\textbf{L'idee.} En faisant courir un pointeur deux fois plus vite qu'un autre, on
obtient deux informations gratuitement : s'ils se rejoignent, il y a un cycle ;
et quand le rapide atteint la fin, le lent est pile au \emph{milieu}. Aucune
memoire supplementaire, un seul passage.

\textbf{Reconnaitre le pattern.} Detection de cycle (liste chainee, suite de
nombres comme \emph{happy number}), recherche du milieu d'une liste, ou detection
d'un nombre duplique en voyant un tableau comme une liste de « pointeurs ».

\begin{infobox}{Quand l'utiliser}
Detecter/localiser un cycle ; trouver le milieu d'une liste ; happy number ;
nombre duplique en \texttt{O(1)} espace.
\end{infobox}

\complexite{O(n)}{O(1)}
\leetcode{Linked List Cycle (LC 141/142), Middle of List (LC 876), Happy Number (LC 202), Find the Duplicate (LC 287).}

\begin{lstlisting}
def has_cycle(head):
    slow = fast = head
    while fast and fast.next:
        slow = slow.next            # +1
        fast = fast.next.next       # +2
        if slow is fast:            # rattrapage -> cycle
            return True
    return False
\end{lstlisting}

\textbf{Piege :} teste \texttt{fast and fast.next} avant \texttt{fast.next.next},
sinon \texttt{NoneType} sur une liste de longueur paire.

% =============================================================================
\section{Merge Intervals (fusion d'intervalles)}
% =============================================================================
On trie les intervalles par debut, puis on fusionne ceux qui se chevauchent au fil
d'un seul parcours.

\textbf{L'idee.} Une fois les intervalles ordonnes par leur debut, deux intervalles
se chevauchent \emph{ssi} le debut du courant est inferieur ou egal a la fin du
dernier conserve. Il suffit alors d'etendre la fin, sinon d'ouvrir un nouvel
intervalle. Le tri transforme un probleme quadratique en un balayage lineaire.

\textbf{Reconnaitre le pattern.} Tout enonce parlant de \emph{plages},
\emph{creneaux}, \emph{reservations} ou \emph{reunions} qui peuvent se recouvrir :
fusionner, inserer, calculer une intersection, ou compter les ressources
simultanees (salles).

\begin{infobox}{Quand l'utiliser}
Fusionner des plages qui se chevauchent ; inserer un intervalle ; intersection de
deux listes ; nombre minimal de salles de reunion.
\end{infobox}

\complexite{O(n log n)}{O(n)}
\leetcode{Merge Intervals (LC 56), Insert Interval (LC 57), Interval Intersection (LC 986), Meeting Rooms II (LC 253).}

\begin{lstlisting}
def merge(intervals):
    intervals.sort(key=lambda p: p[0])       # trier par debut
    res = []
    for start, end in intervals:
        if res and start <= res[-1][1]:      # chevauchement
            res[-1][1] = max(res[-1][1], end)
        else:
            res.append([start, end])
    return res
\end{lstlisting}

\textbf{Piege :} pense au chevauchement « bord a bord » (\texttt{[1,4]} et
\texttt{[4,5]}) : selon l'enonce, \texttt{<=} fusionne, \texttt{<} non.

% =============================================================================
\section{Cyclic Sort (tri cyclique)}
% =============================================================================
Quand un tableau contient des nombres d'une plage connue (\texttt{1..n} ou
\texttt{0..n}), on place chaque valeur a sa position attendue par echanges. Toute
place mal remplie revele un manquant ou un duplique.

\textbf{L'idee.} La valeur \texttt{v} « appartient » a l'index \texttt{v-1}. En
echangeant jusqu'a ce que chaque case soit bien remplie, on trie en \texttt{O(n)}
sans memoire. Ensuite, un simple balayage montre quels index ne correspondent
pas : ce sont les valeurs absentes ou en double.

\textbf{Reconnaitre le pattern.} Tableau de \texttt{n} entiers dans \texttt{[1..n]}
ou \texttt{[0..n]}, et on cherche un \emph{nombre manquant}, un \emph{duplique},
ou \emph{tous les disparus} -- le tout idealement en espace constant.

\begin{infobox}{Quand l'utiliser}
Nombres dans une plage \texttt{[1..n]} ; trouver le manquant, le duplique, ou tous
les disparus en \texttt{O(1)} espace.
\end{infobox}

\complexite{O(n)}{O(1)}
\leetcode{Missing Number (LC 268), Find All Disappeared (LC 448), Set Mismatch (LC 645), First Missing Positive (LC 41).}

\begin{lstlisting}
def cyclic_sort(nums):
    i = 0
    while i < len(nums):
        correct = nums[i] - 1               # place attendue de nums[i]
        if nums[i] != nums[correct]:
            nums[i], nums[correct] = nums[correct], nums[i]
        else:
            i += 1
    return nums
\end{lstlisting}

\textbf{Piege :} n'incremente \texttt{i} que lorsque l'element est deja a sa place,
sinon tu sautes des valeurs mal placees.

% =============================================================================
\section{BFS (parcours en largeur)}
% =============================================================================
On explore le graphe \emph{niveau par niveau} depuis une source, a l'aide d'une
file (queue). Dans un graphe non pondere, BFS donne le \emph{plus court chemin} en
nombre d'aretes.

\textbf{L'idee.} Une file FIFO garantit qu'on visite les noeuds par distance
croissante : d'abord les voisins directs, puis leurs voisins, etc. En traitant un
niveau entier a la fois (boucle sur \texttt{len(queue)}), on sait exactement a
quelle profondeur on se trouve.

\textbf{Reconnaitre le pattern.} « Plus court chemin » / « nombre minimal
d'etapes » dans une grille ou un graphe non pondere, parcours d'arbre par niveaux,
ou propagation simultanee depuis plusieurs sources (multi-source BFS).

\begin{infobox}{Quand l'utiliser}
Plus court chemin non pondere ; parcours par niveaux ; distance minimale dans une
grille ; propagation multi-source.
\end{infobox}

\complexite{O(V + E)}{O(V)}
\leetcode{Binary Tree Level Order (LC 102), Rotting Oranges (LC 994), Shortest Path in Grid (LC 1091).}

\begin{lstlisting}
from collections import deque
def level_order(root):
    if not root: return []
    res, q = [], deque([root])
    while q:
        level = []
        for _ in range(len(q)):        # tout le niveau courant
            node = q.popleft()
            level.append(node.val)
            if node.left:  q.append(node.left)
            if node.right: q.append(node.right)
        res.append(level)
    return res
\end{lstlisting}

\textbf{Piege :} marque un noeud comme visite \emph{au moment de l'enfiler}, pas au
defilement, sinon il peut entrer plusieurs fois dans la file.

% =============================================================================
\section{DFS (parcours en profondeur)}
% =============================================================================
On explore aussi loin que possible le long d'une branche avant de revenir en
arriere (backtrack). Implementation naturelle par recursion (pile d'appels).

\textbf{L'idee.} A partir d'un noeud, on plonge recursivement dans un voisin non
visite, puis le suivant, etc. La pile d'appels memorise le chemin de retour. DFS
epouse parfaitement les structures arborescentes et le decoupage en composantes
connexes.

\textbf{Reconnaitre le pattern.} Compter/explorer des \emph{regions} (iles dans
une grille), enumerer des \emph{chemins}, parcourir un arbre (pre/in/post-order),
ou detecter des composantes connexes.

\begin{infobox}{Quand l'utiliser}
Composantes connexes ; nombre/aire d'iles ; chemins dans un arbre ou graphe ;
exploration exhaustive d'une structure.
\end{infobox}

\complexite{O(V + E)}{O(V) pile}
\leetcode{Number of Islands (LC 200), Max Area of Island (LC 695), Preorder Traversal (LC 144).}

\begin{lstlisting}
def num_islands(grid):
    n, m = len(grid), len(grid[0])
    def sink(r, c):
        if 0 <= r < n and 0 <= c < m and grid[r][c] == "1":
            grid[r][c] = "0"           # marquer visite
            sink(r+1, c); sink(r-1, c); sink(r, c+1); sink(r, c-1)
    count = 0
    for r in range(n):
        for c in range(m):
            if grid[r][c] == "1":
                count += 1; sink(r, c)
    return count
\end{lstlisting}

\textbf{Piege :} sur de grands graphes, la recursion peut depasser la limite de
pile Python -- passe a une pile explicite si besoin.

% =============================================================================
\section{Topological Sort (tri topologique)}
% =============================================================================
Ordonne les sommets d'un graphe oriente \emph{acyclique} (DAG) : pour chaque arete
\texttt{u -> v}, \texttt{u} apparait avant \texttt{v}. Sert a ordonnancer des
taches avec dependances.

\textbf{L'idee (Kahn).} On calcule le \emph{degre entrant} de chaque sommet, puis
on retire d'abord ceux qui n'ont aucune dependance (in-degree 0). Retirer un
sommet decremente le degre de ses successeurs, qui deviennent a leur tour
disponibles. Si l'on ne sort pas tous les sommets, c'est qu'il y a un \emph{cycle}.

\textbf{Reconnaitre le pattern.} Dependances entre elements : ordre de cours
(prerequis), ordre de compilation, planification de taches, resolution de
contraintes « X doit venir avant Y ».

\begin{infobox}{Quand l'utiliser}
Ordonnancer des taches avec prerequis ; detecter un cycle dans un graphe oriente ;
build/compilation ; sequence de cours valide.
\end{infobox}

\complexite{O(V + E)}{O(V)}
\leetcode{Course Schedule (LC 207), Course Schedule II (LC 210), Alien Dictionary (LC 269).}

\begin{lstlisting}
from collections import deque
def topo_order(n, edges):
    adj = {i: [] for i in range(n)}
    indeg = [0]*n
    for u, v in edges:               # u -> v
        adj[u].append(v); indeg[v] += 1
    q = deque(i for i in range(n) if indeg[i] == 0)
    order = []
    while q:
        node = q.popleft(); order.append(node)
        for nxt in adj[node]:
            indeg[nxt] -= 1
            if indeg[nxt] == 0: q.append(nxt)
    return order if len(order) == n else []   # [] => cycle
\end{lstlisting}

\textbf{Piege :} un resultat plus court que \texttt{n} signale un cycle -- toujours
verifier cette condition avant d'utiliser l'ordre.

% =============================================================================
\section{Union-Find / DSU (ensembles disjoints)}
% =============================================================================
Structure qui regroupe des elements en ensembles disjoints avec deux operations
quasi constantes : \texttt{find} (representant) et \texttt{union} (fusion).

\textbf{L'idee.} Chaque ensemble est un arbre pointant vers une racine. \texttt{find}
remonte a la racine ; \texttt{union} accroche une racine sous l'autre. Deux
optimisations rendent le tout quasi \texttt{O(1)} : la \emph{compression de chemin}
(aplatir l'arbre pendant \texttt{find}) et l'\emph{union par rang} (accrocher le
petit arbre sous le grand).

\textbf{Reconnaitre le pattern.} Connexite \emph{dynamique} (aretes ajoutees au
fil de l'eau), comptage de composantes, detection de cycle dans un graphe non
oriente, ou algorithme de Kruskal (arbre couvrant minimal).

\begin{infobox}{Quand l'utiliser}
Composantes connexes dynamiques ; detecter un cycle (graphe non oriente) ;
regrouper des elements ; Kruskal (MST).
\end{infobox}

\complexite{O(alpha(n)) $\approx$ O(1)}{O(n)}
\leetcode{Number of Provinces (LC 547), Redundant Connection (LC 684), Accounts Merge (LC 721).}

\begin{lstlisting}
class UnionFind:
    def __init__(self, n):
        self.parent = list(range(n)); self.rank = [0]*n
    def find(self, x):
        while self.parent[x] != x:
            self.parent[x] = self.parent[self.parent[x]]  # compression
            x = self.parent[x]
        return x
    def union(self, a, b):
        ra, rb = self.find(a), self.find(b)
        if ra == rb: return False        # deja connectes
        if self.rank[ra] < self.rank[rb]: ra, rb = rb, ra
        self.parent[rb] = ra
        if self.rank[ra] == self.rank[rb]: self.rank[ra] += 1
        return True
\end{lstlisting}

\textbf{Piege :} sans compression de chemin \emph{et} union par rang, \texttt{find}
peut degenerer en \texttt{O(n)} sur une chaine.

% =============================================================================
\section{Dynamic Programming (programmation dynamique)}
% =============================================================================
On decompose un probleme en sous-problemes qui \emph{se chevauchent}, et on
memorise leurs solutions pour ne les calculer qu'une seule fois.

\textbf{L'idee.} Quand une solution recursive recalcule sans cesse les memes
sous-problemes, on stocke chaque resultat (memoisation top-down) ou on remplit un
tableau \texttt{dp} des petits cas vers les grands (bottom-up). Deux ingredients
indispensables : \emph{sous-structure optimale} (l'optimum se compose d'optima de
sous-problemes) et \emph{sous-problemes repetes}.

\textbf{Reconnaitre le pattern.} Formulations « nombre de facons de… », « cout
min/max pour… », « est-ce possible… », choix sequentiels avec dependance sur les
etapes precedentes (escaliers, pieces, sac a dos, sous-sequences).

\begin{infobox}{Quand l'utiliser}
« Nombre de facons » ; min/max cout ; possible/impossible ; choix avec
sous-structure optimale et sous-problemes qui se repetent.
\end{infobox}

\complexite{O(etats $\times$ transition)}{O(etats)}
\leetcode{Climbing Stairs (LC 70), Coin Change (LC 322), House Robber (LC 198), LIS (LC 300), Knapsack.}

\begin{lstlisting}
def coin_change(coins, amount):
    dp = [float("inf")] * (amount + 1)
    dp[0] = 0
    for coin in coins:
        for x in range(coin, amount + 1):
            dp[x] = min(dp[x], dp[x - coin] + 1)
    return dp[amount] if dp[amount] != float("inf") else -1
\end{lstlisting}

\textbf{Piege :} definis clairement l'\emph{etat} (que represente \texttt{dp[i]} ?)
et le \emph{cas de base} avant d'ecrire la transition -- 90\% des bugs viennent de
la.

% =============================================================================
\section{Greedy (algorithme glouton)}
% =============================================================================
A chaque etape, on fait le choix \emph{localement optimal} en esperant atteindre
l'optimum global. Rapide, mais valable seulement si le probleme a la « propriete
du choix glouton ».

\textbf{L'idee.} Contrairement a la DP, le glouton ne revient jamais en arriere :
il s'engage sur le meilleur choix immediat. Ca ne fonctionne que si l'on peut
\emph{prouver} qu'un choix local optimal ne ferme pas la porte a l'optimum global
(souvent apres un tri bien choisi).

\textbf{Reconnaitre le pattern.} Ordonnancement d'activites/intervalles,
optimisation avec un critere simple a trier (fin la plus proche, plus gros
d'abord), sauts, monnaie sur systemes canoniques.

\begin{infobox}{Quand l'utiliser}
Choisir/planifier des intervalles ; jump game ; probleme optimisable par un tri +
choix local prouve sur ; alternative rapide a la DP.
\end{infobox}

\complexite{O(n log n) (souvent)}{O(1)}
\leetcode{Jump Game (LC 55/45), Non-overlapping Intervals (LC 435), Candy (LC 135), Best Time to Buy/Sell II (LC 122).}

\begin{lstlisting}
def can_jump(nums):
    reachable = 0
    for i, jump in enumerate(nums):
        if i > reachable:                # trou infranchissable
            return False
        reachable = max(reachable, i + jump)
    return True
\end{lstlisting}

\textbf{Piege :} un glouton « evident » peut etre faux (ex. Coin Change avec des
pieces non canoniques). En cas de doute, verifie avec la DP.

% =============================================================================
\section{Bit Manipulation (manipulation de bits)}
% =============================================================================
On opere directement sur la representation binaire des entiers. Souvent la
solution la plus rapide et en espace constant.

\textbf{L'idee.} Les operateurs bit a bit (\texttt{\&}, \texttt{|}, \texttt{\^},
\texttt{\~}, \texttt{<<}, \texttt{>>}) manipulent tous les bits en parallele.
Quelques identites reviennent sans cesse : \texttt{x \^\ x = 0} (le XOR annule les
doublons), \texttt{x \& (x-1)} efface le bit 1 le plus a droite, \texttt{x \& 1}
teste la parite.

\textbf{Reconnaitre le pattern.} Elements en double sauf un (XOR), comptage de
bits, tests de puissance de deux, ensembles representes en bitmask, ou toute
optimisation memoire/vitesse a base d'entiers.

\begin{infobox}{Quand l'utiliser}
Trouver l'unique non double (XOR) ; compter les bits ; puissance de deux ; masques
de bits (bitmask) pour representer des ensembles.
\end{infobox}

\complexite{O(n) ou O(1)}{O(1)}
\leetcode{Single Number (LC 136), Number of 1 Bits (LC 191), Power of Two (LC 231), Counting Bits (LC 338).}

\begin{lstlisting}
def single_number(nums):
    result = 0
    for x in nums:
        result ^= x            # les doublons s'annulent, reste l'unique
    return result

def count_ones(n):
    count = 0
    while n:
        n &= n - 1             # efface le 1 le plus a droite
        count += 1
    return count
\end{lstlisting}

\textbf{Piege :} en Python les entiers ont une taille illimitee ; pour simuler du
32 bits, masque avec \texttt{\& 0xFFFFFFFF}.

\newpage
% =============================================================================
\section{Fiche recapitulative des complexites}
% =============================================================================
A memoriser avant l'entretien :

\renewcommand{\arraystretch}{1.35}
\begin{longtable}{>{\raggedright\arraybackslash}p{5.2cm}
                  >{\centering\arraybackslash}p{4cm}
                  >{\centering\arraybackslash}p{3cm}}
\rowcolor{primary}
\textbf{\color{white}Pattern} & \textbf{\color{white}Temps} & \textbf{\color{white}Espace} \\
\hline
Hash Map / Set        & \texttt{O(1)} moyen          & \texttt{O(n)} \\
Two Pointers          & \texttt{O(n)}                & \texttt{O(1)} \\
Sliding Window        & \texttt{O(n)}                & \texttt{O(k)} \\
Prefix Sum            & \texttt{O(n)} / \texttt{O(1)}& \texttt{O(n)} \\
Sorting               & \texttt{O(n log n)}          & \texttt{O(n)} \\
Binary Search         & \texttt{O(log n)}            & \texttt{O(1)} \\
Stack                 & \texttt{O(n)}                & \texttt{O(n)} \\
Monotonic Stack       & \texttt{O(n)}                & \texttt{O(n)} \\
Backtracking          & \texttt{O(2\textsuperscript{n})} / \texttt{O(n!)} & \texttt{O(n)} \\
Linked List           & \texttt{O(n)}                & \texttt{O(1)} \\
Fast \& Slow Pointers & \texttt{O(n)}                & \texttt{O(1)} \\
Merge Intervals       & \texttt{O(n log n)}          & \texttt{O(n)} \\
Cyclic Sort           & \texttt{O(n)}                & \texttt{O(1)} \\
BFS                   & \texttt{O(V + E)}            & \texttt{O(V)} \\
DFS                   & \texttt{O(V + E)}            & \texttt{O(V)} \\
Topological Sort      & \texttt{O(V + E)}            & \texttt{O(V)} \\
Union-Find (DSU)      & \texttt{O($\alpha$(n))}      & \texttt{O(n)} \\
Dynamic Programming   & \texttt{O(etats $\times$ tr.)}& \texttt{O(etats)} \\
Greedy                & \texttt{O(n log n)}          & \texttt{O(1)} \\
Bit Manipulation      & \texttt{O(n)} / \texttt{O(1)}& \texttt{O(1)} \\
\end{longtable}

\vspace{0.6cm}
\begin{infobox}{Checklist du jour J}
\begin{itemize}[leftmargin=*,itemsep=2pt,topsep=2pt]
  \item Parler a voix haute -- le raisonnement est note.
  \item Clarifier les contraintes AVANT de coder.
  \item Annoncer la complexite temps + espace.
  \item Traiter les cas limites (vide, 1 element, negatifs, debordement).
  \item Nommer clairement ; code lisible $>$ code astucieux.
  \item Bloque ? Revenir a la force brute, puis optimiser.
  \item Tester a la main avant de dire « c'est bon ».
\end{itemize}
\end{infobox}

\end{document}
